\heading{数学物理方法（下）-模拟题2-期末}

\noteinfo{题目和答案由AI生成。}

\begin{question}
在整条实轴上求解对流--扩散方程初值问题
\[
\begin{cases}
 u_t+v u_x=\kappa u_{xx}, & -\infty<x<\infty,\ t>0,\\
 u(x,0)=e^{-\alpha x^2}, & \alpha>0.
\end{cases}
\]
要求写出显式解。
\begin{solution}

可先作平移变换 $\xi=x-vt$，令 $u(x,t)=w(\xi,t)$，则
\[
 w_t=\kappa w_{\xi\xi},
 \qquad w(\xi,0)=e^{-\alpha \xi^2}.
\]
Gaussian 初值在热方程下保持 Gaussian 形式，故
\[
 w(\xi,t)=\frac{1}{\sqrt{1+4\alpha\kappa t}}
 \exp\left(-\frac{\alpha \xi^2}{1+4\alpha\kappa t}\right).
\]
代回 $\xi=x-vt$，得
\ansmath{u(x,t)=\frac{1}{\sqrt{1+4\alpha\kappa t}}
\exp\left[-\frac{\alpha (x-vt)^2}{1+4\alpha\kappa t}\right].}
\end{solution}
\end{question}

\begin{question}
求三维上半空间
\[
\mathbb R_+^3=\{(x,y,z):z>0\}
\]
中 Dirichlet 问题的 Green 函数：
\[
\nabla^2 G(\mathbf r,\mathbf r')=-\delta(\mathbf r-\mathbf r'),
\qquad G|_{z=0}=0.
\]
并写出位于 $(0,0,h)$ 的单位点源所产生的势函数。
\begin{solution}

用镜像法。记
\[
\mathbf r'=(x',y',z'),
\qquad
\mathbf r'^*=(x',y',-z').
\]
则上半空间 Dirichlet Green 函数为
\ansmath{G(\mathbf r,\mathbf r')=
\frac{1}{4\pi|\mathbf r-\mathbf r'|}-
\frac{1}{4\pi|\mathbf r-\mathbf r'^*|}.}
它在 $z=0$ 上确为零，且满足所需奇性。

若单位点源位于 $(0,0,h)$，则 $\mathbf r_0=(0,0,h)$，其镜像点为 $\mathbf r_0^*=(0,0,-h)$。势函数即为
\ansmath{u(\mathbf r)=\frac{1}{4\pi}\left(\frac{1}{|\mathbf r-\mathbf r_0|}-\frac{1}{|\mathbf r-\mathbf r_0^*|}\right).}
\end{solution}
\end{question}

\begin{question}
考虑 Hilbert 空间 $L^2(0,2\pi)$ 上的算符
\[
\hat L=-i\frac{d}{dx},
\qquad
D(\hat L)=\{u\in H^1(0,2\pi):u(2\pi)=e^{i\theta}u(0)\},
\]
其中 $\theta\in\mathbb R$ 为常数。
\begin{enumerate}[label=（\arabic*）,leftmargin=2.8em,itemsep=0.2em]
\item 证明 $\hat L$ 是自伴的；
\item 求全部本征值与本征函数。
\end{enumerate}
\begin{solution}

\begin{enumerate}[label=（\arabic*）,leftmargin=2.8em,itemsep=0.2em]
\item 对 $u,v\in D(\hat L)$，有
\[
\langle \hat L u,v\rangle-\langle u,\hat L v\rangle
=-i\int_0^{2\pi}u'\bar v\,dx+i\int_0^{2\pi}u\overline{v'}\,dx
=-i[u\bar v]_0^{2\pi}.
\]
而准周期边界给出
\[
u(2\pi)=e^{i\theta}u(0),\qquad v(2\pi)=e^{i\theta}v(0),
\]
故
\[
[u\bar v]_0^{2\pi}=e^{i\theta}u(0)e^{-i\theta}\overline{v(0)}-u(0)\overline{v(0)}=0.
\]
因此 $\hat L$ 对称；其伴随域满足同样边界条件，所以 $\hat L$ 自伴。

\item 本征方程
\[
-i u'=\lambda u
\]
解为
\[
u(x)=C e^{i\lambda x}.
\]
代入边界条件：
\[
e^{i2\pi\lambda}=e^{i\theta}.
\]
故
\[
2\pi\lambda=\theta+2\pi n,
\qquad n\in\mathbb Z.
\]
于是全部本征值为
\ansmath{\lambda_n=n+\frac{\theta}{2\pi},\qquad n\in\mathbb Z,}
相应本征函数可取
\ansmath{u_n(x)=\frac{1}{\sqrt{2\pi}}e^{i(n+\theta/(2\pi))x}.}
\end{enumerate}
\end{solution}
\end{question}

\begin{question}
在区间 $0<x<l$ 上求算符
\[
L=-\frac{d^2}{dx^2}+a^2\qquad (a>0)
\]
配以 Dirichlet 边界条件 $u(0)=u(l)=0$ 的 Green 函数，即求 $G(x,\xi)$ 使得
\[
\left(-\frac{d^2}{dx^2}+a^2\right)G(x,\xi)=\delta(x-\xi),
\qquad G(0,\xi)=G(l,\xi)=0.
\]
并写出方程 $Lu=f(x)$ 的解的积分表示。
\begin{solution}

对 $x\neq \xi$，有齐次方程
\[
-G''+a^2G=0,
\]
其解为双曲函数。为满足左右边界，可取
\[
G(x,\xi)=
\begin{cases}
A(\xi)\sinh(ax), & 0<x<\xi,\\
B(\xi)\sinh(a(l-x)), & \xi<x<l.
\end{cases}
\]
在 $x=\xi$ 处连续，且由方程积分得导数跃变
\[
G_x(\xi+0,\xi)-G_x(\xi-0,\xi)=-1.
\]
解得
\ansmath{G(x,\xi)=\frac{1}{a\sinh(al)}
\sinh(a x_<)\sinh\bigl(a(l-x_>)\bigr),}
其中
\[
x_<=\min\{x,\xi\},\qquad x_>=\max\{x,\xi\}.
\]
因此边值问题
\[
\left(-\frac{d^2}{dx^2}+a^2\right)u=f(x),
\qquad u(0)=u(l)=0
\]
的解为
\ansmath{u(x)=\int_0^l G(x,\xi)f(\xi)\,d\xi.}
\end{solution}
\end{question}


